% Plan: development/outline.md, section 6, item B.2.
% Sources: development/dossiers/dtoc5-optcdeg2.md and its critique,
% development/reviews/sol-dtoc5-exact.md, development/decision-register.md
% (DO-01 to DO-18, N-05, SV-08), development/dossiers/solvers.md (Prop. 16),
% data/numbers.json.
% Round-1 revision (adjudication G5): one certificate box per result, the
% verification-record table merged into the boxes.
\subsection{\inst{dtoc5} and \inst{optcdeg2}}\label{app:dtoc5}

Both instances are Euler transcriptions of optimal control problems with about $5\cdot10^4$ stages, taken from QPLIB and, before that, from CUTEst; both are nonconvex QCQPs with unbounded states, and both certificates are splits (\cref{sec:split-framework}).
For \inst{dtoc5}, \cref{sec:split-affine} proves \cref{prop:dtoc5-identity,thm:dtoc5-bracket}; \cref{app:dtoc5-proof} describes the evaluation and proves \cref{cor:dtoc5-attained}.
For \inst{optcdeg2}, \cref{app:optcdeg2-proof,app:optcdeg2-point} prove \cref{thm:optcdeg2-bound} and the existence of the exactly feasible point it uses.

% ---------------------------------------------------------------------------
\subsubsection{\inst{dtoc5}: stored model}\label{app:dtoc5-model}

The $99{,}999$ continuous variables of the model of \cref{sec:split-affine} are, in the OSIL order, $u_t=x_{t+2}$ ($0\le t<T$) and $y_t=x_{50001+t}$ ($0\le t\le T$).
OSIL row~$t$ reads $-h\,u_t+y_t-y_{t+1}+4h\,y_t^2=0$, with the coefficient strings \texttt{2e-5} for $h$ (also in the objective) and \texttt{8e-5} for $4h$; $y_0=1$ is fixed, and $y_0^2$ contributes the constant $h$ to the objective.
The GAMS file equals that of \inst{QPLIB_8585} apart from comment lines, the formatting of one option and the solve statement (\cref{app:semantics-provenance}); floating-point solves of this model and of the CUTEst problem DTOC5, which has $h$ in place of $4h$, give different optimal values.

All four listed dual bounds are below $0.0025$, against an optimal value of about $5.39$; we take this as a sign that termwise relaxations of the terms $4hy_t^2$ with unbounded $y_t$ are weak, which is an interpretation (\cref{sec:interpretation}), not a tested fact.

% ---------------------------------------------------------------------------
\subsubsection{\inst{dtoc5}: evaluation of the certificate}\label{app:dtoc5-proof}

We use the notation of \cref{prop:dtoc5-identity,thm:dtoc5-bracket}.
The sign convention matters: $c_t$ is the negative of OSIL row~$t$; had the multipliers $\hat\lambda_t$ been attached to the OSIL rows themselves, the coefficient of $y_t^2$ in the Lagrangian would be $h(1+4\hat\lambda_t)$, which is negative wherever $\hat\lambda_t<-\tfrac14$, and the Lagrangian would be unbounded below.

\paragraph{Multipliers and the point.}
Stationarity of the Lagrangian in $u_t$ gives $\lambda_t=-2u_t$, which is why the certificate reads $\hat\lambda$ from the stored exactly feasible point $\hat x$ of \cref{app:points}.
The states $\hat y_1,\dots,\hat y_T$ of $\hat x$ are a refined Newton solution rounded to 30 decimals, and its controls $\hat u_t=50{,}000(\hat y_t-\hat y_{t+1})+4\hat y_t^2$ are defined by the rows, so every row holds exactly.
Every nonterminal control is positive (the smallest is about \sci{1.46}{-5}); the last one, $\hat u_{T-1}\approx-1.96\cdot10^{-26}$, is not used.
Hence $\hat\lambda_t\le0<\tfrac14$ for all $t$, and $A_t\ge h$.

\paragraph{Evaluation.}
Both $f(\hat x)$ and $q(\hat\lambda)$ are rational.
$f(\hat x)$ is computed exactly; it lies below $5.38967211918114046742396649913627186883131341$.
Two codes give rigorous lower bounds on $q(\hat\lambda)$:
\begin{enumerate}
\item[(i)] one reads the GAMS text, whose coefficients equal the OSIL strings exactly, evaluates the polynomial part of $q(\hat\lambda)$ exactly and rounds each of the 49,998 subtracted quotients \emph{up} to the grid $2^{-256}$ before subtracting, which yields a rational $B_{256}\le q(\hat\lambda)$;
\item[(ii)] the other, which has its own OSIL reader, assembles the Lagrangian coefficients from the parsed terms, computes $q(\hat\lambda)$ as an exact reduced rational (its denominator has 9,071,237 bits; the sums use GMP integers), and checks that $f(\hat x)-q(\hat\lambda)$ equals the exactly computed sum of the squares in \eqref{eq:dtoc5-squares}.
\end{enumerate}
Both give $q(\hat\lambda)\ge5.38967211918114046742396649913627186883131268$; the two values differ by about \sci{2.2}{-73}.
The exact gap is $f(\hat x)-q(\hat\lambda)=7.2051\ldots\cdot10^{-43}\le\sci{7.21}{-43}$.
The two 44-decimal displays of \cref{thm:dtoc5-bracket} differ by \sci{7.3}{-43}, more than this gap; \cref{tab:closures} prints 19-decimal displays and therefore marks its gap cell as computed from the certificate ends.
The displayed lower end is a rigorous rational evaluation below the dual function at $\hat\lambda$; we claim neither that $q(\hat\lambda)$ is the optimal value of the Lagrangian dual nor that the duality gap is zero.

\begin{certbox}[Certificate for \cref{thm:dtoc5-bracket}]
\certitem{Inputs}{OSIL file (SHA-256 \texttt{ac7b2d94}); point $\hat x$ (\texttt{6e963788}); exact $f(\hat x)$ (\texttt{8c110402}).}
\certitem{Computation}{all rows and $f(\hat x)$ exactly; $q(\hat\lambda)$ by (i) and (ii); \arith{E} (\texttt{Fraction}, GMP sums).}
\certitem{Data reading}{exact decimals; separately written readers and an exact comparison with the GAMS text confirm every variable, bound, objective term and row.}
\certitem{Implementations}{the exact evaluation (ii) certifies the displayed lower end; the separately written rounded evaluation (i) certifies the same display.}
\certitem{Evidence level}{\evid{hand}+\evid{stored}.}
\certitem{Replay}{Tier~1; about 17\,s for (ii) and 3\,s for (i).}
\end{certbox}

\begin{corollary}[attainment and location]\label{cor:dtoc5-attained}
The optimal value of \inst{dtoc5} is attained.
Every minimizer $x^*$ satisfies $|u^*_t+\hat\lambda_t/2|\le\sci{1.9}{-19}$ for $0\le t\le T-1$ and $|y^*_t-\bar y_t(\hat\lambda)|\le\sci{1.9}{-19}$ for $1\le t\le T-1$.
\end{corollary}

\begin{proof}
Apply \eqref{eq:dtoc5-squares} with $\lambda=\hat\lambda$; here $\hat\lambda_t\le0$, so $A_t\ge h$.
Every feasible $x$ with $f(x)\le f(\hat x)$ has each square term at most $f(\hat x)-q(\hat\lambda)\le\sci{7.21}{-43}$ (\cref{thm:dtoc5-bracket}).
So $u_0,\dots,u_{T-1}$ and $y_1,\dots,y_{T-1}$ are bounded on this level set, and $y_T$ is determined by the last row.
The feasible set is closed and $f$ is continuous, so the level set is compact and nonempty, and $f$ attains its minimum on it.
At a minimizer each deviation is therefore at most $(\sci{7.21}{-43}/h)^{1/2}<\sci{1.9}{-19}$.
\end{proof}

The corollary does not by itself show that the minimizer is unique, and we do not claim uniqueness.

\begin{remark}[Lagrangian and semidefinite relaxations]\label{rem:dtoc5-sdp}
The proof of \cref{prop:dtoc5-identity} collects terms without using the rows, so the identity \eqref{eq:dtoc5-squares} holds with $f(x)+\sum_t\hat\lambda_tc_t(x)$ on the left for every $x$ with $y_0=1$.
The squares vanish at $u_t=-\hat\lambda_t/2$ and $y_t=\bar y_t(\hat\lambda)$, so $q(\hat\lambda)$ is the value of the Lagrangian dual function at $\hat\lambda$, and the optimal value of the Lagrangian dual lies in $[q(\hat\lambda),\vstar]$, an interval of width at most \sci{7.21}{-43}.
The same holds for the order-one semidefinite relaxation, dense (Shor) or sparse.
Each square in \eqref{eq:dtoc5-squares} involves a single variable $z$ and has the form $\gamma(z-\zeta)^2$ with $\gamma>0$.
Its linearization is $\gamma(Z-2\zeta z+\zeta^2)$, where $Z$ stands for $z^2$; this is nonnegative whenever the $2\times2$ moment matrix with entries $1$, $z$, $z$, $Z$ is positive semidefinite, which every moment matrix of either relaxation that contains $z$ ensures.
The linearized rows vanish, so the relaxed objective equals $q(\hat\lambda)$ plus these nonnegative terms.
If $y_0$ is kept as a variable with the row $y_0=1$, a multiple of this row adds the square $A_0(y_0-1)^2$, and the argument is unchanged.
For the source model with coefficient $h$ and 600 to 1000 steps, \citet{waki2006-sums-of-squares-and-semidefinite} observed numerically tight sparse order-one relaxations in floating-point SDP solves (\cref{tab:literature}).
The mechanism is that of Mangasarian-type sufficiency conditions \citep{mangasarian1966-sufficient-conditions-for-the-optimal}: at multipliers taken from a near-optimal point, the Lagrangian is convex.
\end{remark}

The floating-point closure credited in \cref{sec:results-prior} concerns a bounded restriction of \inst{QPLIB_8585}: its MINOTAUR log states that default lower and upper bounds were assumed for 99,983 and 99,997 variables.
Its value $5.3897$ and QPLIB's value $5.38967212$ \citep{furini2026-qplib-official-solution-point-records} agree with ours.

% ---------------------------------------------------------------------------
\subsubsection{\inst{optcdeg2}: stored model}\label{app:optcdeg2-model}

Let $N=50{,}000$ and $h=1/2500$.
The OSIL coefficient strings \texttt{4e-4}, \texttt{8e-6}, \texttt{8e-5} and \texttt{2e-4} are $h$, $a=h/50$, $b=h/5$ and $w=h/2$.
The $150{,}002$ continuous variables are, in the OSIL order, $u_t=x_{2+t}$ for $t\le N-2$, $u_{N-1}=x_{50002}$, $y_t=x_{50003+t}$, $v_t=x_{100004+t}$ for $t\le N-1$, and $v_N=x_{50001}$.
The model of \cref{sec:split-richer} reads
\begin{equation}\label{eq:optcdeg2-model}
\min\ f=w\sum_{t=0}^{N}y_t^2
\quad\text{s.t.}\quad
y_{t+1}=y_t+hv_t,\quad v_{t+1}=v_t+hu_t-ay_t-bv_t^2\quad(t<N),
\end{equation}
with $-\tfrac15\le u_t\le\tfrac15$ (the strings \texttt{-.2} and \texttt{.2}), $v_t\ge-1$ for $1\le t\le N-1$, $y_0=10$ and $v_0=v_N=0$; the positions $y_1,\dots,y_N$ are free, and the velocities have no upper bound.
The objective is convex; the only nonconvex terms are $-bv_t^2$ in the velocity rows.
The GAMS file equals that of \inst{QPLIB_8803} apart from comment lines, a redundant declaration, the formatting of one option and the solve statement (\cref{app:semantics-provenance}).

Gurobi's listed dual bound $292.41713458$ equals, to the displayed digits, the objective of the listed point p2, which MINLPLib records with violation $10^{-6}$.
Evaluated under exact semantics, p2 violates 46,687 rows by more than $10^{-8}$, by at most \sci{9.97}{-7}.
The listed bound is valid: it lies at least $1.45$ below the certified bound of \cref{thm:optcdeg2-bound}.
The value of p2 is attained by no exactly feasible point (category~A, \cref{tab:claims}).

% ---------------------------------------------------------------------------
\subsubsection{\inst{optcdeg2}: the quadratic split}\label{app:optcdeg2-proof}

The certificate uses a split whose function on the state $(y_t,v_t)$ has the form
\begin{equation}\label{eq:optcdeg2-S}
\varphi_{t-1}(y,v)=p^y_t\,y+p^v_t\,v+\frac{q_t}2\,(v-\bar v_t)^2,\qquad 1\le t\le N,
\end{equation}
with rational data $(p^y_t,p^v_t,q_t,\bar v_t)$; as in \cref{sec:split-framework}, the index of a split function is that of the stage to whose right it acts, so $\varphi_{t-1}$ links the stages $t-1$ and $t$.
Write $U=[-\tfrac15,\tfrac15]$ and $g(v)=v-bv^2$.

\begin{lemma}[split bound for \inst{optcdeg2}]\label{lem:optcdeg2-krotov}
Let $\varphi_0,\dots,\varphi_{N-1}\colon\R^2\to\R$ be arbitrary, and for $1\le t\le N-1$ let $V_t\subseteq\R$ contain $v_t$ for every exactly feasible point.
Put $\rho_t(y,v,u)=wy^2+\varphi_t(y+hv,\ g(v)+hu-ay)-\varphi_{t-1}(y,v)$.
Then every exactly feasible point satisfies $f\ge m_0+\sum_{t=1}^{N-1}m_t+m_N$, where
\begin{gather*}
m_0=\inf_{u\in U}\bigl[100w+\varphi_0(10,\,hu-10a)\bigr],\qquad
m_N=\inf_{y\in\R}\bigl[wy^2-\varphi_{N-1}(y,0)\bigr],\\
m_t=\inf\bigl\{\rho_t(y,v,u):\ y\in\R,\ v\in V_t,\ u\in U\bigr\}\quad(1\le t\le N-1).
\end{gather*}
\end{lemma}

\begin{proof}
For an exactly feasible point, $(y_{t+1},v_{t+1})$ is the image of $(y_t,v_t,u_t)$ under the dynamics, so the split terms telescope:
\[
f=\bigl[wy_0^2+\varphi_0(y_1,v_1)\bigr]+\sum_{t=1}^{N-1}\rho_t(y_t,v_t,u_t)+\bigl[wy_N^2-\varphi_{N-1}(y_N,v_N)\bigr].
\]
Here $y_0=10$ and $v_0=0$, so $(y_1,v_1)=(10,\,hu_0-10a)$, and $u_t\in U$, $v_t\in V_t$ and $v_N=0$.
Each bracket is at least its infimum.
\end{proof}

This is \cref{lem:split-bound}(a) for the stages $(y_t,v_t,u_t)$, with the derived enclosures $\R\times V_t\times U$ (\cref{rem:split-enclosures}).

\begin{lemma}[velocity enclosure]\label{lem:optcdeg2-enclosure}
Define intervals by a forward pass, $Y_0=\{10\}$, $V_0=\{0\}$,
\[
Y_{t+1}=Y_t+hV_t,\qquad V_{t+1}=\bigl(g(V_t)+hU-aY_t\bigr)\cap[-1,\infty)\quad(t+1\le N-1),\qquad V_N=\{0\},
\]
followed by a backward pass for $t=N-1,\dots,1$,
\[
Y_t\leftarrow Y_t\cap(Y_{t+1}-hV_t),\qquad V_t\leftarrow V_t\cap g^{-1}\bigl(V_{t+1}-hU+aY_t\bigr),
\]
where $g^{-1}$ is the inverse of $g$ on $(-\infty,1/(2b))$.
Suppose that $\sup V_t<1/(2b)=6250$ for all $t$ after the forward pass, and that every step is rounded outward.
Then $(y_t,v_t)\in Y_t\times V_t$ for every exactly feasible point and every $t$.
\end{lemma}

\begin{proof}
By induction.
In the forward pass, if $(y_t,v_t)\in Y_t\times V_t$, then $y_{t+1}=y_t+hv_t\in Y_t+hV_t$ and $v_{t+1}=g(v_t)+hu_t-ay_t\in g(V_t)+hU-aY_t$, and $v_{t+1}\ge-1$ for $t+1\le N-1$; since $g$ is increasing on $(-\infty,6250)$, $g(V_t)$ is the interval between the images of the endpoints.
The row $v_N=0$ gives $V_N=\{0\}$.
In the backward pass, by induction from $t+1$, $y_t=y_{t+1}-hv_t$ and $g(v_t)=v_{t+1}-hu_t+ay_t$, and $v_t$ lies on the increasing branch, so $v_t\in g^{-1}(V_{t+1}-hU+aY_t)$.
Outward rounding only enlarges the intervals.
\end{proof}

The computation of these intervals is exact: each step is evaluated in rational arithmetic and rounded outward to the dyadic grid $2^{-110}$; $g^{-1}(s)$ is enclosed by rational points $x^-\le x^+$ for which $g(x^-)\le s\le g(x^+)$ is checked exactly; and the final bounds are rounded outward to binary64.
The forward bounds stay below $1.22$, and the computation also checks that $0\in g(V_{N-1})+hU-aY_{N-1}$, which every feasible point requires.
The result satisfies $V_t\subseteq[-1,\,1.2135]$, with widths at most $2.09$.

\begin{lemma}[exact stage reduction]\label{lem:optcdeg2-stage}
Fix $1\le t\le N-1$, let $\varphi_{t-1},\varphi_t$ be of the form \eqref{eq:optcdeg2-S}, and put
\begin{gather*}
A=w+\tfrac12q_{t+1}a^2,\qquad B=p^y_{t+1}-p^y_t-a\,p^v_{t+1},\\
K_2=\frac{q_{t+1}w}{2A},\qquad K_1=p^v_{t+1}+\frac{a\,q_{t+1}B}{2A},\qquad K_0=-\frac{B^2}{4A},\\
\Gamma(v)=\bigl(h\,p^y_{t+1}-p^v_t\bigr)v+p^v_{t+1}\bar v_{t+1}+K_0-\frac{q_t}2(v-\bar v_t)^2 .
\end{gather*}
If $A>0$, then $\min_{y\in\R}\rho_t(y,v,u)=G(v,e):=\Gamma(v)+K_2e^2+K_1e$ with $e=g(v)+hu-\bar v_{t+1}$.
For $u\in U$, $e$ ranges over $[e_-(v),e_+(v)]=g(v)-\bar v_{t+1}+[-h/5,\,h/5]$, and
\[
m_t\ \ge\ \min\Bigl\{\inf_{v\in V_t}G\bigl(v,e_-(v)\bigr),\ \inf_{v\in V_t}G\bigl(v,e_+(v)\bigr),\ \gamma\Bigr\},
\]
where $\gamma=+\infty$ if $K_2\le0$, and otherwise $\gamma=\inf_{v\in R^*}\bigl[\Gamma(v)-K_1^2/(4K_2)\bigr]$ for any set $R^*\supseteq\{v\in V_t:\ e_-(v)\le-K_1/(2K_2)\le e_+(v)\}$ (with $\gamma=+\infty$ if $R^*=\emptyset$).
Equality holds if $R^*$ is exactly this set.
\end{lemma}

\begin{proof}
Put $y'=y+hv$ and $v'=g(v)+hu-ay$, so that $v'-\bar v_{t+1}=e-ay$.
Expanding \eqref{eq:optcdeg2-S},
\[
\rho_t=Ay^2+\bigl(B-q_{t+1}a\,e\bigr)y+h\,p^y_{t+1}v+p^v_{t+1}(e+\bar v_{t+1})+\frac{q_{t+1}}2e^2-p^v_tv-\frac{q_t}2(v-\bar v_t)^2 .
\]
For $A>0$, minimizing over $y\in\R$ subtracts $(B-q_{t+1}ae)^2/(4A)$.
The coefficient of $e^2$ becomes $q_{t+1}/2-q_{t+1}^2a^2/(4A)=q_{t+1}(2A-q_{t+1}a^2)/(4A)=K_2$, the coefficient of $e$ becomes $K_1$, and the remaining terms form $\Gamma(v)$.
For fixed $v$, the map $u\mapsto e$ is affine and increasing, with range $[e_-(v),e_+(v)]$.
A quadratic in $e$ attains its minimum over an interval at an endpoint, or, if $K_2>0$, possibly at the vertex $-K_1/(2K_2)$, where its value is $\Gamma(v)-K_1^2/(4K_2)$; this value is a lower bound on $G(v,\cdot)$ for every $v$, so enlarging $R^*$ can only lower $\gamma$.
\end{proof}

Since $g$ is quadratic, $G(v,e_\pm(v))$ is a polynomial of degree at most $4$ in $v$.

\begin{lemma}[rigorous univariate minimization]\label{lem:optcdeg2-quartic}
Let $G$ be a nonconstant real polynomial and $\alpha<\beta$ with $G'(\alpha)G'(\beta)\neq0$.
Let $[l_k,r_k]\subset(\alpha,\beta)$, $k=1,\dots,K$, satisfy $G'(l_k)<0<G'(r_k)$, and suppose that every point of $(\alpha,\beta)$ at which $G'$ changes sign from negative to positive lies in one of these intervals.
If $M_k\ge\max_{[l_k,r_k]}|G''|$, then
\[
\min_{[\alpha,\beta]}G\ \ge\ \min\Bigl\{G(\alpha),\ G(\beta),\ \min_k\bigl[\min\bigl(G(l_k),G(r_k)\bigr)-M_k(r_k-l_k)^2\bigr]\Bigr\}.
\]
\end{lemma}

\begin{proof}
Let $x_0$ minimize $G$ over $[\alpha,\beta]$; if $x_0$ is an endpoint, there is nothing to prove.
Otherwise $G'(x_0)=0$.
Since $G'$ is a nonzero polynomial, it has finitely many zeros, so it has a constant sign on each side of $x_0$ near $x_0$.
It must be negative on the left and positive on the right, since otherwise $G$ would be smaller at a nearby point.
So $x_0\in[l_k,r_k]$ for some $k$, and Taylor's theorem gives $G(l_k)=G(x_0)+\tfrac12G''(\xi)(l_k-x_0)^2\le G(x_0)+\tfrac12M_k(r_k-l_k)^2$.
\end{proof}


\paragraph{Certificate data.}
The data $(p^y_t,p^v_t,q_t,\bar v_t)$ are binary64 numbers stored with the certificate and read as exact rationals; any data give a valid bound.
Here $\bar v_t$ is the velocity of a refined floating-point KKT point, and $p^y_t,p^v_t$ are its discrete costates.
With $q\equiv0$, \cref{lem:optcdeg2-krotov} gives the Lagrangian bound of \cref{prop:split-affine}, whose stage residual contains the term $-b\,p^v_{t+1}v^2$.
This term is concave in $v$ on the two arcs where the certified control is $-\tfrac15$ ($t<3091$ and $t>47290$), because $p^v_{t+1}>0$ there ($p^v_1\approx124.09$, $p^v_N\approx0.595$), and there the affine split is not minimized at the trajectory (\cref{sec:split-richer}); with an earlier set of multipliers, the certified affine bound is $293.2500703$, about $0.63$ below the optimum.

The quadratic term adds curvature in $v$ only.
On the two arcs with control $-\tfrac15$, $q_t$ is chosen so that the $v$-curvature of $\rho_t$ at the trajectory, $-2b\,p^v_{t+1}+q_{t+1}(1-2b\bar v_t)^2-q_t$, equals $2\eta h$, with margin $\eta=1$ on the first arc and $\eta=0.05$ on the last.
The recursion runs backward from $q_{3091}=0$ on the first arc, reaching about $-34.1$ at the start, and forward from $q_{47291}=0$ on the last, reaching about $0.286$ at $t=N$.
On the middle arc $q\equiv0$, because there $p^v<0$ and the affine residual is already convex in $v$.
Since $v_0$ and $v_N$ are fixed, the curvature of the split at the two ends never enters the bound, so the schedules can start from $q=0$ at each switch and run outward; and $q=0$ at both switching stages keeps the two stages with fractional control nearly exact.
The validity of the bound does not depend on these design choices.
\Cref{fig:optcdeg2-calibration} shows the control of the certified trajectory, the velocity costate and the curvature schedule.

\begin{figure}[tb]
\centering
\includegraphics[width=0.85\textwidth]{fig-optcdeg2-calibration}
\caption{Split data of the \inst{optcdeg2} certificate against time $th$.
(a)~Control $u_t$ of the certified trajectory.
(b)~Velocity costate $p^v_t$; it is positive on both arcs with $u=-\frac15$.
(c)~Curvature $q_t$ of the split functions; it is zero on the middle arc and at both switching stages (dashed lines, $t=3091$ and $t=47290$).
Panels (b) and (c) use a symmetric logarithmic scale, linear near~0.
Data: the stored binary64 split data.}
\label{fig:optcdeg2-calibration}
\end{figure}

\begin{proof}[Proof of \cref{thm:optcdeg2-bound}]
Apply \cref{lem:optcdeg2-krotov} with the functions \eqref{eq:optcdeg2-S} and the intervals $V_t$ of \cref{lem:optcdeg2-enclosure}.
For $1\le t\le N-1$, \cref{lem:optcdeg2-stage} applies because $A>0$ at every stage, which the computation checks exactly.
It reduces $m_t$ to the minimization of two polynomials of degree at most $4$ over $V_t$; the third term never arises, because at every stage either $K_2\le0$ or a computed superset of the vertex set is empty.
Polynomials of degree at most $2$ are minimized in closed form.
For the 11,596 quartics, an exact Sturm sequence of $G'$ isolates its distinct real roots in $(\alpha,\beta)$, after an exact check that $G'(\alpha)G'(\beta)\ne0$.
For each root at which $G'$ changes sign from negative to positive, a bracket with $G'(l)<0<G'(r)$ is checked exactly, with floating-point roots used only as seeds (7,785 brackets; no bisection fallback and no root of even multiplicity was needed).
\Cref{lem:optcdeg2-quartic}, with $M_k$ the sum of the absolute coefficients of $G''$ times powers of $\max(|l_k|,|r_k|)$ (in the shifted variable $v-\bar v_t$), then gives a rational $\tilde m_t\le m_t$; the largest slack $M_k(r_k-l_k)^2$ is \sci{9.5}{-29}.
The end terms are explicit: $m_0$ is the minimum of a quadratic in $u$ over $U$ (about $1014.99647$), and $m_N=-(p^y_N)^2/(4w)-\tfrac12q_N\bar v_N^2=-(p^y_N)^2/(2h)$ because $\bar v_N=0$ (about $-1.86\cdot10^{-4}$).
Rounding $m_0$, every $\tilde m_t$ and $m_N$ down to the grid $2^{-200}$ and summing exactly gives a rational $B\le\vstar$ with $B\ge293.87607509587509237940$, the rigorous evaluation in \cref{thm:optcdeg2-bound}; its truncated display is $293.87607509587509237$.
All arithmetic is in Python integers and \texttt{Fraction}.
With the upper bound of \cref{prop:optcdeg2-point}, $\gapabs\le\sci{8.98}{-16}$ and $\gaprel\le\sci{3.06}{-18}$.
\end{proof}

The computation also checks the reduction of \cref{lem:optcdeg2-stage} against the definition of $\rho_t$, exactly, at three points $(v,u)$ in each of 115 stages, among them the seven stages around each switch.

Two further checks support the exact code and are not part of the proof.
A separately written sampling search over 197 stages found no point below the computed stage minimum (best values within \sci{2.8}{-29} of it).
Along the certified trajectory rounded to $2^{-300}$ at each step (near-feasible, not exactly feasible), every stage residual is at least its computed minimum, and the differences sum to $8.978\cdot10^{-16}$, the gap between $B$ and the upper bound of \cref{prop:optcdeg2-point}; stage 3091 accounts for about \sci{8.9777}{-16}, stage 47290 for about \sci{4.08}{-20} and every other stage for at most \sci{7.2}{-25}, the two stages with fractional control, where $p^v_{t+1}$ is small but not zero (\sci{7.6}{-12} and $-\sci{9.7}{-16}$).
The verification principle is classical (\cref{sec:split-framework}); the model-specific function and its exact check are the contribution.

\begin{certbox}[Certificate for \cref{thm:optcdeg2-bound}]
\certitem{Inputs}{OSIL file (SHA-256 \texttt{6a9bcd91}); split data (\texttt{08311d33}).}
\certitem{Computation}{$V_t$ on the grid $2^{-110}$; 11,596 quartic pieces by exact Sturm isolation; floors on the grid $2^{-200}$ summed exactly; \arith{E}.}
\certitem{Data reading}{exact decimals, asserted by the reader; a separately written reader and an exact comparison with the GAMS text confirm every row, bound and objective term.}
\certitem{Implementations}{the second code (exact) certifies the displayed bound; the first code, in binary64 intervals (\arith{F}, 162 geometric cells per stage, refined by bisection), certifies the weaker $293.8760750958728$, each stage bound at least \sci{3.6}{-19} below the exact one.}
\certitem{Evidence level}{\evid{stored}.}
\certitem{Replay}{Tier~1; 6\,s and 19\,s (enclosure and stage minima; 16 to 18\,s for the latter in other runs); 8\,s for the interval code.}
\end{certbox}

% ---------------------------------------------------------------------------
\subsubsection{\inst{optcdeg2}: the exactly feasible point}\label{app:optcdeg2-point}

\begin{proposition}[an exactly feasible point]\label{prop:optcdeg2-point}
Let $u_t=-\tfrac15$ for $t\le3090$ and $t\ge47291$, $u_t=\tfrac15$ for $3092\le t\le47289$, $u_{3091}=1506828432880435/2^{55}\approx0.041823$, and let $u_{47290}=\sigma$ be free.
Let $y_t(\sigma),v_t(\sigma)$ be defined by the rows of \eqref{eq:optcdeg2-model} from $y_0=10$, $v_0=0$.
For a rational $\tilde\sigma\approx0.0951284142971001$, there is $\sigma^*\in[\tilde\sigma-10^{-60},\tilde\sigma+10^{-60}]$ such that the resulting point is exactly feasible, and its objective is at most $293.87607509587509328$.
\end{proposition}

\begin{proof}
For every $\sigma$, all rows hold by construction, and $y_0=10$, $v_0=0$.
All controls lie in $U$ for $\sigma$ in the bracket.
The function $\sigma\mapsto v_N(\sigma)$ is a polynomial, hence continuous.
An interval simulation of the rows in integer arithmetic, with all quantities scaled by $2^{320}$ and rounded by floor and ceiling, gives $v_N<0$ at the left end of the bracket and $v_N>0$ at the right end (about $\mp3.8\cdot10^{-64}$).
Over the whole bracket it gives $v_t(\sigma)\ge-0.8316$ for $1\le t\le N-1$ and encloses the objective in an interval of width below \sci{4.8}{-64} whose upper end is $293.8760750958750932772142114247080159340995594910\ldots$.
By the intermediate value theorem, some $\sigma^*$ in the bracket gives $v_N=0$; the point is then exactly feasible, and its objective lies in the enclosure.
\end{proof}

The control of the certified point is $\pm\tfrac15$ except at stages 3091 and 47290; we claim this structure only for the certified point, not for every minimizer.
Our floating-point KKT point violates rows by up to \sci{8.9}{-16}, and its objective $293.87607509588105$ lies at least \sci{5.95}{-12} above the upper bound of \cref{prop:optcdeg2-point}; this says nothing about attainability.

\begin{certbox}[Certificate for \cref{prop:optcdeg2-point}]
\certitem{Inputs}{OSIL file (SHA-256 \texttt{6a9bcd91}); control vector (binary64, \texttt{8206c29d}).}
\certitem{Computation}{interval simulation of all rows over the bracket on the grid $2^{-320}$; sign change of $v_N$; bounds; objective enclosure; \arith{E}.}
\certitem{Data reading}{exact decimals, asserted as rationals.}
\certitem{Implementations}{the integer code certifies the displayed upper bound, which therefore does not depend on \texttt{mpmath}; a separately written \texttt{mpmath} code (\arith{I}, 200 bits, bracket of half-width $10^{-40}$) certifies the same display.}
\certitem{Evidence level}{\evid{stored}.}
\certitem{Replay}{Tier~1; 1.5\,s (integer code) and 9\,s (\texttt{mpmath} code).}
\end{certbox}

% ---------------------------------------------------------------------------
\subsubsection{\inst{optcdeg2}: MINOTAUR's report}\label{app:optcdeg2-minotaur}

\Cref{sec:solvers-invalid} proves \cref{prop:minotaur-optcdeg2} from \cref{prop:optcdeg2-point}, and \cref{app:solvers-published} compares the files and the log.
According to the log, the report occurs in the first presolve, before MINOTAUR's quadratic handler assigns default bounds to unbounded variables, so such default bounds play no role; the cause of the report is unknown, and we attribute no defect to MINOTAUR.
